\section*{Chapter \ref{ch:b3:origin-of-elements} --- Where the Elements Come From}
% Values from the Book 5 scratch calculator calc/ch30.py.

\begin{solution}{exo:b3:origin-of-elements:1}
(a)~$\mathrm{Si/Fe} = 10^{7.51 - 7.46} = 10^{0.05} = 1.12$ and
$\mathrm{O/Fe} = 10^{8.69 - 7.46} = 10^{1.23} = 17.0$: about one silicon
atom and seventeen oxygen atoms per iron atom. (b)~$\mathrm{Fe/Au} =
10^{7.46 - 0.91} = 10^{6.55} = \num{3.55e6}$ and $\mathrm{Pb/Au} =
10^{1.95 - 0.91} = 10^{1.04} = 11.0$: beyond the iron peak the abundances
fall by millions, but lead, on its own peak, is eleven times more
abundant than gold.
\end{solution}

\begin{solution}{exo:b3:origin-of-elements:2}
(a)~$10^{7.46 - 12} = \num{2.884e-5}$, that is \num{28.8} iron atoms per
million hydrogen atoms. (b)~A hundred times fewer: \num{0.288} per
million, one iron atom per $1/\num{2.884e-7} = \num{3.47e6}$ hydrogen
atoms. (c)~$[\mathrm{Fe/H}] = \log_{10}(\num{e-7}/\num{2.884e-5}) =
-2.46$.
\end{solution}

\begin{solution}{exo:b3:origin-of-elements:3}
$^{124}\mathrm{Xe}$ and $^{126}\mathrm{Xe}$ lie on the proton-rich side,
off the s path, and are shielded by the stable $^{124}\mathrm{Te}$ and
$^{126}\mathrm{Te}$: \omterm{def:b3:origin-of-elements:p-process}{p nuclei}. $^{128}\mathrm{Xe}$ and $^{130}\mathrm{Xe}$
are on the s path but shielded from the \omterm{def:b3:origin-of-elements:s-process}{r process} by the stable
$^{128}\mathrm{Te}$ and $^{130}\mathrm{Te}$: s-only. $^{129}\mathrm{Xe}$,
$^{131}\mathrm{Xe}$ and $^{132}\mathrm{Xe}$ are on the s path and have no
stable isobar of lower $Z$ (iodine has only 127): both processes.
$^{134}\mathrm{Xe}$ and $^{136}\mathrm{Xe}$ lie beyond
$^{133}\mathrm{Xe}$, where the s path leaves xenon, and are not
shielded: r-only.
\end{solution}

\begin{solution}{exo:b3:origin-of-elements:4}
(a)~$t = d/c = 40\times\qty{3.086e22}{m}/\qty{2.998e8}{m.s^{-1}} =
\qty{4.12e15}{s}$, about \num{130} million years. (b)~$\Delta v/c \le
1.7/\num{4.12e15} = \num{4.1e-16}$. This is an upper bound, since part or
all of the delay may come from the emission of the gamma rays after the
merger.
\end{solution}

\begin{solution}{exo:b3:origin-of-elements:5}
(a)~$295\times0.1067 = \qty{31.48}{mb}$, $832\times0.0373 =
\qty{31.03}{mb}$, $155\times0.1984 = \qty{30.75}{mb}$. (b)~From
$^{122}\mathrm{Te}$: $N_{123} = 31.48/832 = \num{0.03783}$, \qty{1.4}{\percent}
above the measured value, and $N_{124} = 31.48/155 = \num{0.2031}$,
\qty{2.4}{\percent} above. Abundances that differ by a factor of five are
predicted to within a few per cent from the cross-sections alone.
(c)~$^{128}\mathrm{Te}$ is r-only: the s flow leaves tellurium at
$^{127}\mathrm{Te}$, which decays, so no steady s flow passes through
$^{128}\mathrm{Te}$ and its abundance was set by the \omterm{def:b3:origin-of-elements:s-process}{r process}.
\end{solution}

\begin{solution}{exo:b3:origin-of-elements:6}
(a)~$\tau_0\langle\sigma\rangle = 0.306\times155 = 47.4$, so
$f = 1/(1 + 1/47.4) = \num{0.979}$ for $^{124}\mathrm{Te}$; for
$^{138}\mathrm{Ba}$, $0.306\times4.0 = 1.224$ and $f = 1/(1 + 0.817) =
\num{0.550}$: the magic nucleus passes on only \qty{55}{\percent} of the
flow. (b)~The six factors are \num{0.959}, \num{0.550}, \num{0.908},
\num{0.771}, \num{0.972} and \num{0.915}; their product is \num{0.328}:
$\sigma N$ falls by a factor of \num{3.05}. (c)~The measured fall is
$27.6/7.72 = 3.57$. The model gets the size of the step right to within
\qty{15}{\percent}, nearly all of it from $^{138}\mathrm{Ba}$ and
$^{140}\mathrm{Ce}$, the two nuclei with $N = 82$.
\end{solution}

\begin{solution}{exo:b3:origin-of-elements:7}
(a)~$kT = 8\times10^3\times\qty{1.602e-19}{J} = \qty{1.282e-15}{J}$, and
$v_T = \sqrt{2\times\num{1.282e-15}/\num{1.6749e-27}} =
\qty{1.24e6}{m.s^{-1}}$. (b)~With $\langle\sigma\rangle =
\qty{155e-31}{m^2}$, $\lambda = n_\mathrm{n}\langle\sigma\rangle v_T$
is \qty{1.92e-11}{s^{-1}} for $n_\mathrm{n} = \qty{e12}{m^{-3}}$
($10^6$ per cubic centimetre): a mean time of \qty{5.2e10}{s}, about
\num{1650} years; then \num{16.5} years at $10^8$ and \num{0.165}~year,
two months, at \qty{e10}{cm^{-3}}. (c)~$^{127}\mathrm{Te}$ decays in
hours, far faster than any of these capture times: it always decays and
the path does not split there. $^{134}\mathrm{Cs}$ lives about two
years, between the capture times at $10^8$ and \qty{e10}{cm^{-3}}: at
the high densities of the helium flashes a good part of it captures
before it decays. It is a branching point, and the abundances on either
side of it measure the neutron density.
\end{solution}

\begin{solution}{exo:b3:origin-of-elements:8}
(a)~$10^{6.30 - 3.14} = 1445$ times rarer than calcium and
$10^{4.97 - 3.14} = 68$ times rarer than titanium. (b)~$10^{8.46 - 1.38}
= \num{1.2e7}$. (c)~Scandium shows the odd--even effect: nuclei with
even numbers of protons are more tightly bound and are built by the
capture of helium nuclei in the burning stages of massive stars, which
favour them. Beryllium shows the trough of lithium, beryllium and boron:
these fragile nuclei are destroyed at the temperatures of stellar
interiors, and helium burning goes from helium straight to carbon,
without making them.
\end{solution}

\begin{solution}{exo:b3:origin-of-elements:9}
(a)~Let $\alpha$ be the \omterm{def:b3:origin-of-elements:alpha}{alpha elements} and $\mathrm{Fe}_\mathrm{cc}$,
$\mathrm{Fe}_\mathrm{Ia}$ the iron from the two kinds of supernova. On
the plateau, $\alpha/\mathrm{Fe}_\mathrm{cc} = 10^{P}$ times the solar
ratio. At $[\mathrm{Fe/H}] = 0$, $\alpha/(\mathrm{Fe}_\mathrm{cc} +
\mathrm{Fe}_\mathrm{Ia})$ is solar. Dividing,
$\mathrm{Fe}_\mathrm{cc}/(\mathrm{Fe}_\mathrm{cc} + \mathrm{Fe}_\mathrm{Ia})
= 10^{-P}$, and the Type~Ia share is $1 - 10^{-P} = 1 - 10^{-0.4} =
\num{0.602}$. (b)~$10^{-P} = 1/3$ gives $P = \log_{10}3 = 0.48$.
(c)~The Type~Ia supernovae begin after the same delay in both galaxies.
When they began in the dwarf galaxy, its gas had reached only
$[\mathrm{Fe/H}] = -2$, against $-1$ in the Galaxy: it had turned its gas
into stars ten times more slowly.
\end{solution}

\begin{solution}{exo:b3:origin-of-elements:10}
(a)~For zirconium,
\[
\frac{(X/\mathrm{Mg})_\mathrm{Earth}}{(X/\mathrm{Mg})_\mathrm{CI}} =
\frac{10.5/228000}{3.96/95870} = 1.11 ;
\]
likewise $(2670/228000)/(5010/95870) = 0.224$ for sodium,
\num{0.190} for potassium and \num{0.118} for rubidium. (b)~Zirconium, the most refractory, stands
close to its chondritic proportion. The three moderately volatile
alkali metals, condensing at \qtyrange{800}{1006}{K}, are depleted by
factors of five to nine, rubidium, the most volatile, the most. The
silicate Earth follows the \omterm{def:b3:origin-of-elements:volatility}{volatility trend}. (c)~Magnesium is a major,
refractory element that condenses with the bulk of the rock, and it is
lithophile: it stayed in the mantle when the core formed. A siderophile
reference such as iron or nickel would mix the loss to the core with the
loss by volatility.
\end{solution}

\begin{solution}{exo:b3:origin-of-elements:11}
(a)~When a mass $\mathrm{d}s$ is locked into stars, the wind removes
$\lambda\,\mathrm{d}s$ of gas of \omterm{def:b3:origin-of-elements:metallicity}{metallicity} $Z$:
$\mathrm{d}M_\mathrm{g} = -(1 + \lambda)\,\mathrm{d}s$ and
$\mathrm{d}(ZM_\mathrm{g}) = y\,\mathrm{d}s - Z\,\mathrm{d}s -
\lambda Z\,\mathrm{d}s$. Expanding the second and subtracting
$Z\,\mathrm{d}M_\mathrm{g}$ leaves $M_\mathrm{g}\,\mathrm{d}Z =
y\,\mathrm{d}s = -\frac{y}{1 + \lambda}\,\mathrm{d}M_\mathrm{g}$, which
integrates to the result. (b)~$\ln(1/0.29) = 1.238$: $Z = 0.0155$ for
$\lambda = 0$, and half as much, $0.0077$, for $\lambda = 1$. (c)~A wind
lowers the effective yield by the factor $1 + \lambda$. In a galaxy of
small mass, supernovae drive gas out more easily, $\lambda$ is large,
and the metals are lost with the gas.
\end{solution}

\begin{solution}{exo:b3:origin-of-elements:12}
(a)~$N_\mathrm{Li}/N_\mathrm{H} = 10^{0.96 - 12} = \num{9.1e-12}$ in the
Sun and $10^{3.25 - 12} = \num{1.78e-9}$ in the CI chondrites. The Sun
keeps $10^{-2.29} = \num{0.0051}$ of its lithium: it has destroyed
\qty{99.5}{\percent} of it. (b)~$\num{1.78e-9}/\num{1.45e-10} = 12.3$ and
$\num{1.78e-9}/\num{4.72e-10} = 3.8$. (c)~The Solar System holds almost
four times more lithium than the Big Bang made: most of its lithium was
made later, by some other process. The oldest stars hold three times
less than the Big Bang made: either they have destroyed part of it, as
the Sun has, or the prediction misses something; the two comparisons
alone do not say which.
\end{solution}

\begin{solution}{pb:b3:origin-of-elements:1}

\textbf{Part I.}

\textbf{1.} $10^{12 - 7.46} = 10^{4.54} = \num{34700}$.

\textbf{2.} $N_\mathrm{He}/N_\mathrm{H} = 10^{10.914 - 12} = \num{0.0820}$, so
$4N_\mathrm{He}/N_\mathrm{H} = \num{0.328}$, against $Y/X =
0.2423/0.7438 = \num{0.326}$. They agree to \qty{0.7}{\percent}; the
small gap is there because a helium atom weighs a little less than four
hydrogen atoms.

\textbf{3.} $\log_{10}\bigl((1/350000)/\num{2.884e-5}\bigr) = -5.544 + 4.540 =
-1.00$.

\textbf{4.} $[\mathrm{O/H}] = -1.00 + 0.40 = -0.60$, so
$\log\varepsilon(\mathrm{O}) = 8.69 - 0.60 = 8.09$. The star has
$10^{8.09 - 6.46} = 42.7$ oxygen atoms per iron atom, against
$10^{1.23} = 17.0$ in the Sun: $10^{0.4} = 2.5$ times more.

\textbf{5.} $Z \approx 0.0139\times10^{-0.60} = \num{0.0035}$, a quarter of the
Sun's.

\textbf{6.} Iron has a great many lines in the spectrum of a cool star, while
oxygen has few and faint ones: iron is measured far more easily and
precisely, and stands in for the metals.

\textbf{Part II.}

\textbf{7.} Each sits behind a stable isobar of lower $Z$:
$^{122}\mathrm{Sn}$, $^{123}\mathrm{Sb}$ and $^{124}\mathrm{Sn}$ stop
the r-process decay chains before they reach tellurium. All three lie
on the s path, which runs through tellurium from $^{122}\mathrm{Te}$ to
$^{126}\mathrm{Te}$.

\textbf{8.} \num{31.48}, \num{31.03} and \qty{30.75}{mb}; mean
\qty{31.09}{mb}.

\textbf{9.} The mean of the other two is $(31.48 + 30.75)/2 = \qty{31.11}{mb}$,
so $N_{123} = 31.11/832 = \num{0.03740}$, only \qty{0.3}{\percent}
above the measured \num{0.0373}.

\textbf{10.} $1.329\times44.4 = \qty{59.0}{mb}$, nearly twice the plateau.
$^{128}\mathrm{Te}$ is r-only: the s path leaves tellurium at
$^{127}\mathrm{Te}$, so its abundance was made by the \omterm{def:b3:origin-of-elements:s-process}{r process} and
owes nothing to the s flow.

\textbf{11.} $1 - 7.43/8.53 = \num{0.13}$: \qty{13}{\percent} of the flow went
round $^{148}\mathrm{Sm}$, by captures on the unstable
$^{147}\mathrm{Nd}$, $^{147}\mathrm{Pm}$ and $^{148}\mathrm{Pm}$ before
they could decay to samarium.

\textbf{12.} With a closed shell of 82 neutrons, these nuclei bind an extra
neutron weakly and have very small capture cross-sections (\qty{4.0}{mb}
for $^{138}\mathrm{Ba}$). The flow piles up on them and only part of it
passes on: they are bottlenecks.

\textbf{Part III.}

\textbf{13.} $\mu_1 = 13.7/(13.7 + 33.4) = 13.7/47.1 = \num{0.291}$.

\textbf{14.} $\ln(1/0.291) = \num{1.235}$, so $y = 0.0154/1.235 = \num{0.0125}$,
\num{0.81} times $Z_\odot$.

\textbf{15.} At $\mu = 0.9$, $\ln(1/0.9) = 0.1054$, so $Z = 0.0125\times0.1054 =
\num{0.00131}$. Since $0.00131/0.0154 = 0.0851$, $[\mathrm{Fe/H}] =
-1.07$. At $\mu = 0.5$, $Z = 0.0125\times0.693 = \num{0.00864}$ and
$[\mathrm{Fe/H}] = -0.25$.

\textbf{16.} $Z = 0.00154$, so $\mu = \mathrm{e}^{-0.00154/0.01247} =
\mathrm{e}^{-0.1235} = \num{0.884}$.

\textbf{17.} $1 - 0.884 = \num{0.116}$: \qty{12}{\percent} of the mass.

\textbf{18.} $y = 0.0154/\ln 10 = \num{0.0067}$, about half as much.

\textbf{Part IV.}

\textbf{19.} The mass of stars with \omterm{def:b3:origin-of-elements:metallicity}{metallicity} below $Z$ is
$M(1 - \mathrm{e}^{-Z/y})$, so per unit $Z$ it is
$(M/y)\,\mathrm{e}^{-Z/y}$. Since $Z = Z_\odot10^{[\mathrm{Fe/H}]}$,
$\mathrm{d}Z = Z\ln 10\,\mathrm{d}[\mathrm{Fe/H}]$, and per unit
$[\mathrm{Fe/H}]$ the distribution is proportional to
$Z\mathrm{e}^{-Z/y}$. Its derivative,
$(1 - Z/y)\mathrm{e}^{-Z/y}$, vanishes at $Z = y$: the peak is at
$[\mathrm{Fe/H}] = \log_{10}0.810 = -0.09$.

\textbf{20.} $\num{0.116}/(1 - 0.291) = 0.1162/0.7091 = \num{0.164}$.

\textbf{21.} $Z/y = 1.235\times10^{-0.5} = 0.3905$, $1 - \mathrm{e}^{-0.3905} =
0.3233$, and $0.3233/0.7091 = \num{0.456}$.

\textbf{22.} $6670\times0.1638 = 1093$, against 26: forty times too many.

\textbf{23.} $26/902 = \qty{2.9}{\percent}$, still more than five times below
the box. The model's stars are the long-lived ones, those of at most
about a solar mass, which have all survived since they formed; among
the more massive F stars, the oldest, which are also the most
metal-poor, have died, and the full sample under-counts them.

\textbf{24.} $(1 - 0.1^{0.1})/(1 - 0.1) = 0.2057/0.9 = \num{0.23}$: worse.
Raising the gas fraction does not help either: as $\mu_1$ tends to 1
the fraction tends to $0.1$, still more than three times the
\qty{2.9}{\percent} observed. No gas fraction rescues the closed box.

\textbf{25.} The closed box predicts that
\textbf{\qty[text-series-to-math=true]{16}{\percent}} of the long-lived
stars of the solar neighbourhood have $[\mathrm{Fe/H}] < -1$, against
\qty{3}{\percent} at most observed. The disc was not a closed box: it
grew by a slow infall of metal-poor gas, so that little gas ever sat
waiting at low \omterm{def:b3:origin-of-elements:metallicity}{metallicity} to form stars.
\end{solution}
